\section{Síntesis y ampliación}

El valor del overview de V5 está en desplegar el Transformer, que suele presentarse como un solo architecture diagram, en su ciclo de vida completo. Embedding/attention definen la interfaz de cómputo; la estructura de los datos y el schedule determinan los gradientes; CoT y search aumentan el test-time compute; la preference optimization modifica el comportamiento a largo plazo; el agent coloca al modelo en un ciclo cerrado con el entorno; la visión y la fMRI muestran la transferencia de la token interface; y el continual learning se pregunta cómo cambia realmente un modelo después de su despliegue.

\subsection{Doce conclusiones centrales}

\begin{enumerate}
  \item El Transformer block es un componente básico de los sistemas de IA modernos, no la explicación de un producto completo.
  \item Q/K/V cumplen, respectivamente, los papeles de consulta, índice de coincidencia y contenido que se agrega.
  \item La calidad, la estructura, la diversidad, la escala y el schedule de los datos deben diseñarse en conjunto.
  \item La comparación con el aprendizaje humano puede generar hipótesis experimentales, pero no demuestra por sí misma mecanismos similares a los del cerebro.
  \item El mixed result de TinyDialogues indica que una diverse mixture suele ser más robusta que datos elegidos por una sola intuición.
  \item La ganancia del preentrenamiento en dos fases depende del blend y de la duration; el sobremuestreo produce rendimientos decrecientes.
  \item CoT, tree, program y decomposition modifican la inference computation; no son una misma forma de «razonamiento».
  \item La diferencia central entre los preference methods está en la fuente de supervisión, el baseline y las objective assumptions.
  \item El self-improvement de un agent requiere verification externa, permisos y condiciones de parada.
  \item La transferencia a ViT/VLM/fMRI depende, antes que nada, de una token interface y un objective correctos.
  \item Los scaling limits deben juzgarse por la pendiente local y los nuevos cuellos de botella, no con consignas que proclamen su fin o su continuidad ilimitada.
  \item El continual learning debe adquirir nuevas capacidades, controlar el olvido e indicar dónde ocurre la actualización, todo a la vez.
\end{enumerate}

\subsection{Una tabla que articula toda la clase}

\begin{table}[H]
\centering
\small
\begin{tabular}{@{}L{0.17\textwidth}L{0.20\textwidth}L{0.28\textwidth}L{0.25\textwidth}@{}}
\toprule
Nivel & Objeto principal & Métodos típicos & Pregunta de aceptación \\
\midrule
Architecture & token interaction & attention, ViT, cross-attention & ¿La representación y la complejidad son razonables? \\
Data & distribución de entrenamiento & TinyDialogues, two-phase blend & ¿Son estables las múltiples métricas, la tasa de repetición y la scale? \\
Inference & test-time compute & CoT, ToT, tools & Presupuesto, faithfulness, verificación \\
Preference & comportamiento de los parámetros & RLHF, DPO, GRPO, KTO & judge, reward, diferencias entre grupos \\
Agent & ciclo cerrado con el entorno & ReAct, reflection, LATS & Permisos, outcome, recuperación \\
Adaptation & actualización tras el despliegue & editing, MoE, continual FT & plasticity, forgetting, rollback \\
\bottomrule
\end{tabular}
\caption{Mapa de seis niveles de un sistema Transformer, del cómputo a la adaptación continua.}
\end{table}

\subsection{Lista de verificación para la investigación práctica}

\begin{importantbox}{Antes de iniciar un nuevo proyecto con Transformer, responda}
\begin{enumerate}
  \item ¿Cómo se tokeniza la entrada y qué relaciones requieren attention?
  \item ¿Qué señal de los datos de entrenamiento supervisa la capacidad objetivo?
  \item ¿La mejora proviene de más parámetros, más tokens, un schedule distinto o test-time search?
  \item ¿Se usa human/AI feedback? ¿Quién audita al judge?
  \item ¿Cuáles son los límites de permisos de las herramientas, la memory y el environment?
  \item ¿La evaluación cubre distribution shift, costo, calibration y recuperación ante fallas?
  \item ¿La actualización tras el despliegue es memory externa, editing local o aprendizaje de pesos?
  \item ¿La nueva capacidad se obtiene a costa de catastrophic forgetting?
\end{enumerate}
\end{importantbox}

\subsection{Conjunto mínimo de registros para la reproducción}

Ya sea que se reproduzca un método de data, reasoning, preference, agent o continual learning, deben guardarse source/version, el hash de los datos, las semillas aleatorias, el modelo y el tokenizer, el optimizer/schedule, el compute total de entrenamiento e inferencia, el protocol de evaluación completo, las muestras fallidas y las reglas de juicio humano. Si se usa un AI judge, herramientas o una base de recuperación, también hay que fijar su versión y sus permisos. Solo con toda esta información otro equipo puede reproducir, atribuir y auditar las mejoras que muestran las figuras; de lo contrario, incluso la curva más atractiva de un overview sirve solo como pista y no como fundamento para decisiones de ingeniería.

\subsection{Preguntas de autoevaluación}

\begin{enumerate}
  \item ¿Cuáles son las formas matriciales y la semántica de Q, K y V? ¿Por qué se divide entre $\sqrt{d_k}$?
  \item ¿En qué se diferencian el static embedding y la contextual representation?
  \item ¿Por qué TinyDialogues no respalda la conclusión simple de que «el corpus puramente infantil es el mejor»?
  \item ¿Por qué en el preentrenamiento en dos fases el blend y la duration deben someterse a ablación por separado?
  \item ¿Qué recurso agrega cada uno de CoT, ToT, PoT y el computation graph?
  \item ¿En qué difieren RLHF, DPO, RLAIF, GRPO y KTO en cuanto a supervisión e infraestructura?
  \item ¿Por qué la self-refinement puede reforzar errores cuando no hay un verifier externo?
  \item ¿Por qué ViT, CLIP y VLM no son sinónimos?
  \item ¿Qué aprende el masked objective de un fMRI foundation model y qué no puede demostrar?
  \item ¿Qué diferencia hay entre una caída de la pendiente local del scaling y que «la scaling law deje de cumplirse»?
  \item ¿Qué estado actualizan, respectivamente, RAG, model editing y continual fine-tuning?
  \item ¿Cómo medir a la vez la plasticity y el forgetting en el continual learning?
\end{enumerate}

\subsection{Lecturas adicionales}

\begin{itemize}
  \item Vaswani et al., \emph{Attention Is All You Need}: Q/K/V, multi-head attention y encoder-decoder.
  \item Dosovitskiy et al., \emph{An Image is Worth 16x16 Words}: Vision Transformer.
  \item Wei et al., \emph{Chain-of-Thought Prompting Elicits Reasoning in Large Language Models}.
  \item Rafailov et al., \emph{Direct Preference Optimization}: preference pairs y policy objective.
  \item Yao et al., \emph{ReAct} y \emph{Tree of Thoughts}: reasoning/action y search.
  \item Meng et al., \emph{ROME}; Mitchell et al., \emph{MEMIT}: model editing.
  \item Los primary papers de TinyDialogues y del preentrenamiento en dos fases citados en las Lecture slides: dos tipos de experimentos sobre la estructura de los datos y el schedule.
\end{itemize}

\end{document}
